\section{Cómo medir la inteligencia: supervivencia, exploración, automatización y economía de los token}

La sección anterior trazó una hoja de ruta; esta sección responde a la pregunta insistente de Zhang Xiaojun (张小珺): «¿cuál es la esencia de la inteligencia?». Guangmi (广密) no dio una definición rigurosa, sino que describió la evolución de la inteligencia humana con tres palabras clave: supervivencia, exploración y automatización. Aplicado a los modelos, el progreso de la inteligencia puede entenderse como la capacidad del modelo de adaptarse mejor a su entorno, descubrir nuevas estrategias, completar tareas de forma automática y realizar trabajos más complejos con un costo en token menor o más eficiente.

\begin{figure}[H]
\centering
\includegraphics[width=0.92\textwidth]{intelligence-measurement.png}
\caption{Cómo se mide el progreso de la inteligencia: de la supervivencia, la exploración y la automatización al consumo de token y al ciclo cerrado de tareas. Diagrama conceptual propio, elaborado a partir de la conversación de 00:43:00--00:48:03.}
\end{figure}

\begin{knowledgebox}{Lectura de la figura: la inteligencia no es una cifra de benchmark}
En la figura aparecen supervivencia, exploración, automatización, Token, tarea y retroalimentación. Al leerla conviene notar lo siguiente: un benchmark es un corte de medición, mientras que la inteligencia real se parece más a un ciclo cerrado de tareas. Que el modelo pueda explorar de forma independiente, elegir herramientas, consumir una cantidad razonable de token y completar tareas verificables es una forma de medición más cercana a la era de los Agent.
\end{knowledgebox}

\subsection{Por qué importa el consumo de token}

Esta sección explica la economía de los token. En la entrevista se menciona que una conversación ordinaria con un chatbot puede consumir algunos miles de token, que una tarea de búsqueda o de investigación puede consumir más, y que una tarea promedio de Manus puede llegar a cientos de miles de token. Las cifras no deben tomarse como estadísticas precisas; lo importante es la tendencia: las tareas de un Agent consumen más razonamiento, contexto, llamadas a herramientas y estados intermedios que las tareas de conversación. Si la fijación de precios de un producto sigue copiando la cuota mensual del SaaS, puede quedar desalineada respecto del costo real de cómputo.

\begin{importantbox}{Experiencia práctica: el precio de un producto de IA debe considerar el costo en token}
El SaaS tradicional tiene un costo marginal bajo, mientras que un Agent puede consumir en cada tarea una gran cantidad de token, de llamadas a herramientas y de cómputo externo. Si el producto se vende con una cuota mensual fija, pero los usuarios ejecutan muchas tareas de alto costo, la economía de los token terminará volviéndose en contra del modelo de negocio.
\end{importantbox}

\subsection{De long context a long-term memory}

Esta sección prepara el terreno para la discusión posterior sobre los Agent. Guangmi menciona que el siguiente milestone de la AGI podría ser la long-term memory, que sustituiría a la simple long context. La long context consiste en meter más material en la ventana actual; la long-term memory es la capacidad del modelo de acumular, seleccionar, actualizar y recuperar de forma continua la experiencia previa. La primera resuelve «ver más en esta ronda»; la segunda, «entender más y hacerlo mejor cuanto más tiempo se usa».

\begin{knowledgebox}{Términos clave: Long Context y Long-term Memory}
\noindent\begin{tabularx}{\textwidth}{p{0.22\textwidth}p{0.34\textwidth}X}
\toprule
Concepto & Significado & Función en un Agent \\
\midrule
Long Context & La ventana de una sola inferencia admite material más extenso & Permite que el modelo lea de una vez más archivos, páginas web e historial. \\
Long-term Memory & Guardar y actualizar experiencia entre sesiones y entre tareas & Permite que el modelo forme un estado persistente del usuario, del proyecto y del entorno. \\
Task State & Objetivo, pasos, errores y productos de la tarea actual & Determina si el Agent puede ejecutar tareas de largo alcance. \\
Experience Reuse & Transferir éxitos y fracasos pasados a tareas nuevas & Es una de las capacidades previas al Online Learning. \\
\bottomrule
\end{tabularx}
\end{knowledgebox}

\subsection{Resumen de la sección}

La medición de la inteligencia no puede basarse solo en la puntuación de preguntas aisladas. En la era de los Agent importan más el ciclo cerrado de tareas, la retroalimentación del entorno, el costo en token, la memoria de largo plazo y la reutilización de la experiencia. En conjunto, estas variables determinan si el modelo realmente pasa de «saber responder» a «saber hacer».
